\section{Conclusion}
\label{sec:conclusion}

We built a benchmark for daily weather forecasting on tropical-monsoon station observations, with the controls
needed to interpret a tropical result: a same-network temperate reference, a reanalysis track at identical
coordinates, and a gap-matched context control applied to models and baselines alike. We evaluated 8 zero-shot
time-series foundation models against 10 trained and statistical baselines over 170{,}571 rolling-origin windows,
treating the forecast origin as the unit of resampling.

The picture is mixed. Foundation models are the strongest family, but on station observations most of their
advantage is not statistically separable. Nothing beats a day-of-year climatology on tropical rainfall beyond a
one-day lead; reanalysis flatters every model relative to the stations it represents; and skill reverses between
monsoon break and active spells. The long-lead tropical temperature deficit that half the roster shows is mostly,
though not entirely, a consequence of sparser station context rather than of the tropics.

For practitioners in data-sparse regions the practical reading is twofold. Zero-shot models are already
affordable and already competitive on smooth variables; for rainfall, a well-built climatology remains the thing
to beat. For benchmark builders the reading is that scoring target, regime stratification and context availability
are not implementation details: each one changed a conclusion in this study.

Two extensions would sharpen the picture. Few-shot adaptation is cheap for several models in the roster and would
test whether the gaps we report are properties of zero-shot use or of the models themselves. Running the same
protocol on a second tropical network would show whether the long-lead temperature deficit is Bangladeshi or
tropical. The harness, splits and cached scores are released so that either is a single run against unchanged
windows.
